\documentclass[]{article}
\usepackage{amsfonts, amssymb, amsmath}
\usepackage{float}
% To turn off indenting (Paragraphs)
\parindent 0px

\usepackage{enumerate}
\begin{document}
	
	The distributive property states that $a (b+c) = ab + bc + ac$, for all $a, b, c, \in \mathbb{R}$ \\[6pt]
	$A$ can accept $[1, 2, 3]$ but not ${1, 2, 3}$ so we do $\{1, 2, 3\}$
	
	It does not accepts \$. So we put a slash before.\\[6pt]
	 
	 $$2(\frac{1}{x^2 - 1})$$. So we do: 
	 $$2\left (\frac{1}{x^2 - 1}\right) $$
	$$2\left [\frac{1}{x^2 - 1}\right ] $$ $$2\left \{\frac{1}{x^2 - 1}\right \} $$ $$2\left | \frac{1}{x^2 - 1}\right | $$
	
	$$\left.\frac{dy}{dx}\right|_{x=1} $$
	
	Tables: 
	
	\begin{tabular}{|c||c|c|c|c|c|} \hline
		$x$ & 1 & 2 & 3 & 4 & 5 \\ \hline
		$f(x)$ & 10 & 20 & 30 & 40 & 50 \\ \hline
	\end{tabular}
	
	\vspace{0.66cm}
	\begin{table}[H] 
	\centering % for data to be centered
	
	% for spacing
	\def \arraystretch{1}
	% If a column has paragraph , it doesnot fit so we have to {c|c|p{3in}} , and it tells about width of data 
	\begin{tabular}{|c||c|c|c|c|c|} \hline
		$x$ & 11 & 22 & 33 & 44 & 55 \\ \hline
		$f(x)$ & 110 & 220 & 330 & 440 & 550 \\ \hline
	\end{tabular}
	\caption{This is caption of the data }
	\end{table}
	% To be code in correct destination we have to use flow package
	
	
	
	Arrays: 
	% everything in alaign is in math mode 
	
	% To alaign with = put & before =
	\begin{align}
		5x^2+3 &= 0 \text{ This is text}\\ 
		a+3 &= b^2-3d \\
		d^2-d-5 = 0
	\end{align}
	% If we want to remove line number, use {align*}
	
	\begin{enumerate} \setcounter{enumi}{6}
		% This allows to be started after 6 
		\item Ali
		\item Ahmad
		\begin{enumerate}
			\item bin
			\item Safeer
		\end{enumerate}
	\end{enumerate}
	Bulleted List
		\begin{itemize}
		\item Ali
		\item Ahmad
		\item Bin Safeer
	\end{itemize}
	
	% for custom names before bullets, we have to import enumerate
	\vspace{2in}
	Changed: 
	\begin{enumerate}[A.]
		\item Ali
		\item Ahmad
		\begin{enumerate}
			\item[] bin  % add[] to remove numbering
			\item Safeer
		\end{enumerate}
	\end{enumerate}
	
\end{document}